\section{模型分析与评价}

\subsection{模型的整体结构}

三个问题共用同一套太阳位置、镜面跟踪、光线反射和损失计算模型，区别只在于待求变量逐步增加。问题一在给定镜场上完成光学性能评价，为后续优化提供统一的目标函数与约束判定器；问题二在统一镜面规格条件下联合搜索塔位、布局和镜面尺寸；问题三进一步开放逐镜宽度、高度和安装高度，并以面积加权方式计算异质镜场功率。这样的递进结构保证了三问数值能够在同一物理口径下比较，也避免在优化阶段使用与评价阶段不一致的近似指标。

\subsection{模型优点}

首先，模型保持了从太阳矢量到接收器命中的完整几何链条。每面定日镜均依据入射方向和指向接收器中心的方向实时确定法向量，余弦损失、阴影遮挡、大气透射和截断损失均在同一光线样本上计算，因此各效率分量与最终功率之间具有明确对应关系。对于问题三的异质镜面，采用镜面面积加权而非逐镜等权，使效率汇总能够严格还原功率计算式。

其次，优化过程把容量约束置于单位面积功率比较之前。问题二和问题三均先要求年平均输出热功率不低于60 MW，再最大化单位镜面面积年平均输出热功率，从而排除了面积很小、效率较高但不能完成供热任务的伪优方案。场界、禁布区、异尺寸间距、宽高关系及离地净空均在候选生成和正式复核阶段再次检查，使所得方案同时满足光学目标和几何可实施条件。

最后，模型采用分层精度控制搜索误差。低分辨率用于大范围筛选，高分辨率和多随机种子只用于候选复核。问题一的分辨率与随机种子检验表明评价结果较为稳定；问题二、问题三的最终方案也均经过两级分辨率和四种子检验。问题三细化过程中还主动剔除了在高分辨率下跌破60 MW的候选，说明容量判定以正式复核结果为准，而不是以搜索阶段的乐观估计为准。

\subsection{模型局限与改进方向}

本文结论适用于题目给定的60个离散太阳状态及所建立的理想化光学模型。太阳盘采用均匀pillbox角分布，尚未引入周日亮度、镜面斜率误差、跟踪误差和长期积灰；接收器按几何命中统计截断效率，也未进一步模拟接收器表面热流密度与热损失。因此，所得功率表示规定状态下的光学输出，不应直接解释为包含天气和设备可用率的实际电站年发电量。

此外，问题二和问题三含有离散镜数、非光滑遮挡关系及大量连续变量，粒子群、遗传算法与差分进化给出的是记录搜索范围和预算下的最佳已验证可行解，而非数学意义上的全局最优解。特别是两问的容量裕量均较小，工程设计时还需为制造、跟踪和气象不确定性预留更大的安全裕度。后续可引入实测DNI权重、非均匀太阳形状和误差卷积，并把接收器峰值热流及容量裕量纳入稳健多目标优化，以提高工程适用性。
